\chapter*{Acknowledgements}
\addcontentsline{toc}{chapter}{\protect\numberline{}Acknowledgements}

Over the years, I have been asked a number of times whether I would recommend starting a PhD program. My answer has always been a resounding "No.\llap{"} followed by a brief pause and "\llap{.}..unless you can find a really good lab"\llap{.} This answer is usually followed by a follow-up question: "What makes a good lab?" First and foremost, the staff and students should be friendly and intelligent. You will draw on them for inspiration and companionship, and the effects of good colleagues on your own work will be profound. Second, there must be infrastructure and funding available to support good research. Both of these are easy for the would-be student to verify: go and visit, and take a look. Finally, the person in charge should be a visionary. He or she \textit{must} have a clear and unique idea of what the future will look like, and how the lab will contribute to, and interact with it. Only then can truly novel and original output be produced consistently.

I was fortunate to have ended up, mostly by chance, in a lab where all three of these conditions were met exceedingly. At Aerts lab, I have met the kindest and most intelligent people who made research both fun and inspiring. During (and after) my tenure there, funding was ample - enabling us to invest in crazy, low-percentage ideas. Admittedly, most of the time we were just figuring out and doing things for the first time ourselves. A lab where this is possible can in practice only exist under good leadership. Staying at the leading edge requires the courage to commit fully to unique ideas, and to effectively bet your continued existence on risky hypotheses over and over again. Keeping all members of an organization aligned during the stressful times that these projects invariably generate is challenging in itself, let alone succeeding in the science to boot.

I would like to thank my advisor Stein, for building this environment. It is thanks to you that I find myself in the midst of an area of science that walks hand in hand with the most incredible revolution of our time. I would also like to thank my supervisor Suresh for his excellent daily support. We endured so many long nights in the lab building HyDrop together! Third, I want to thank Jasper for his supervision during my master thesis, which led to this work. All three of you have left an imprint on who I am today.

I want to thank Linda, for supporting me throughout. I have learned a lot from you.

Finally, I want to thank my parents, Dominique \& Johan, for giving me everything.